\clearpage
\section{Core proof details, counterexamples and verification scope}
\label{app:core-proof-details}

The main text states the hypotheses, conclusions and proof ideas of the
MM/HC spectral identification and the fixed-exponent correction model.
This appendix supplies the complete arguments and their limits. It is a
direct reference target, not a background process. No guarded citation,
command name or compiler result supplies an omitted mathematical premise.

\subsection{Complete proofs of the conditional spectral statements}
\label{app:hc-spectral-proof}

\begin{proof}[Proof of Lemma~\ref{lem:diagonal-traversal-scalar-action}]
Let $S_C=\mathrm{im}(C)$ in the bounded complex Hilbert realization of
the lemma. Commutation implies invariance: if $x=Cx$, then
$C\mathcal Mx=\mathcal MCx=\mathcal Mx$.
By the complete-eigenvector hypothesis, every $x\in S_C$ is a limit
of finite linear combinations of eigenvectors in $S_C$. The
single-eigenvalue hypothesis assigns eigenvalue $\lambda_0$ to each
of them. Thus $\mathcal Mx=\lambda_0 x$ on that span, and boundedness
extends the identity to its closure, all of $S_C$. In finite dimensions
this is exactly diagonalizability together with one eigenvalue.
\end{proof}

\begin{proof}[Proof of Theorem~\ref{thm:hc-spectral-projector}]
Set $\mathcal S_0=\ker(\mathcal M-\lambda_0 I)$. The preceding
lemma gives $S_C\subseteq\mathcal S_0$. If $x\in\mathcal S_0$,
commutation gives
\[
 (\mathcal M-\lambda_0 I)(I-C)x
   =(I-C)(\mathcal M-\lambda_0 I)x=0,
 \qquad C(I-C)x=0.
\]
The complement-exclusion hypothesis makes $(I-C)x=0$, so $x=Cx$
and $\mathcal S_0\subseteq S_C$. Hence the two spaces are equal.
The assumed orthogonal pure-point decomposition defines the spectral
projector onto the full $\lambda_0$ eigenspace as its orthogonal
projection. Since $C$ is orthogonal with exactly that image, it is
this projector. Neither the decomposition nor complement exclusion
is inferred from $[\mathcal M,C]=0$ alone.
\end{proof}

The scalar compatibility in Corollary~\ref{cor:M-HC-fixpoint}
follows by applying the scalar action to $Cx$ and then using
commutation. A native identification with closure-stable states is
separate. Unbounded realizations need a common invariant domain and
an appropriate spectral-measure argument; no such extension is proved here.

\subsection{Independent spectral failure modes}
\label{app:hc-spectral-counterexamples}

\paragraph{Several eigenvalues without Jordan chains.}
For $\mathcal M=\mathrm{diag}(1,2,3)$ and
$C=\mathrm{diag}(1,1,0)$, both operators are diagonal and commute.
The compensated image nevertheless carries eigenvalues $1$ and $2$.
There are no nontrivial Jordan blocks, but no scalar action on this image.
The cascade statement $\hat K(L_n)=L_{n-1}$ concerns another operator
and does not establish the Jordan or spectral structure of $\mathcal M$.

\paragraph{One eigenvalue without diagonalizability.}
The block
\[
 J_2=\begin{pmatrix}2&1\\0&2\end{pmatrix}
\]
has only eigenvalue $2$ but is not $2I$. The complete-eigenvector
condition cannot be dropped merely because the eigenvalue is unique.

\paragraph{Scalar action without full spectral selection.}
For $\mathcal M=\mathrm{diag}(2,2,3)$ and
$C=\mathrm{diag}(1,0,0)$, the identities
$\mathcal MC=C\mathcal M=2C$ hold, but the full eigenvalue-$2$
projector is $\mathrm{diag}(1,1,0)$. This distinguishes image
scalarity from exclusion of equal eigenvectors on the complement.

\paragraph{A projection does not determine the full operator.}
The same $C=\mathrm{diag}(1,1,0)$ obeys the same scalar compatibility
with either $\mathrm{diag}(2,2,3)$ or $\mathrm{diag}(2,2,4)$.
The action on $\ker C$ remains free until independently specified.

The older scalar-action argument invoked irreducibility not supplied
by Lemma~\ref{lem:closure-forcing-lemma}. A possible finite-dimensional
repair would specify an irreducible represented algebra on $S_C$
commuting with $\mathcal M$. Every $\mathcal M$ eigenspace is then
invariant under that algebra. Over the complex field an eigenvector
exists, and irreducibility forces its eigenspace to be all of $S_C$.
This would establish scalarity, but the represented algebra, its
irreducibility and its commutation must actually be exhibited.

\subsection{Complete proof for the fixed-exponent ansatz}
\label{app:fixed-exponent-proof}

\begin{proof}[Proof of Theorem~\ref{thm:self-consistent-gL}]
Equation (SC2) supplies the stated relation between $\mu$ and
$\bsym$. Equation (SC1) is assumed as a correction ansatz; it is
not the Weyl-density ratio rejected by the earlier numerical test.
For $\mu>1$ both sides are positive. Taking logarithms reduces
(SC1)--(SC2) to
\[
 f(\mu)=\tfrac12\log3+\tfrac12\log\frac{\mu-1}{\mu+1}
                 +\frac{a}{\mu^5}=0,
 \qquad a=\frac{\gamma_1}{4\pi}.
\]
Differentiation gives
\[
 f'(\mu)=\frac1{\mu^2-1}-\frac{5a}{\mu^6}.
\]
For $t=\mu^2>1$, the exact factorization
\[
 t^3-\frac{27}{4}(t-1)
       =\left(t-\frac32\right)^2(t+3)\geq0
\]
proves $\mu^6/(\mu^2-1)\geq27/4$. The hypothesis
$0<a<27/20$ therefore makes $f'>0$ everywhere on $(1,\infty)$.
Moreover $f(1+)=-\infty$ and $f(2)=a/32>0$. The intermediate
value theorem gives a root in $(1,2)$ and strict monotonicity gives
uniqueness on $(1,\infty)$. This establishes the scalar model theorem.
\end{proof}

At the fixed numeric input $\gamma_1=14.134725141734694\ldots$,
$a\simeq1.12480568777618$ lies in the stated interval. Bisection
gives $\mu^*\simeq1.863662249889$ and
$\bsym\simeq0.549175395418$. The approximately $0.03957\%$
difference from the inherited comparison value $1.8644$ is arithmetic
proximity, not an observational uncertainty or a derivation of the ansatz.

\begin{center}
\small
\begin{tabular}{@{}p{0.24\textwidth}p{0.37\textwidth}p{0.27\textwidth}@{}}
\toprule
Comparison & Declared condition & Values $(\gL,\bsym)$\\
\midrule
Continuum & $\gamma_1\to0$ in (SC1) & $(2,1/\sqrt3)$\\
Self-referential & $\gL=1/\bsym$ together with (SC2) &
$(\mathbb T,1/\mathbb T)$, $\mathbb T\approx1.839$\\
Fixed exponent & First-zero input, exponent $5$ assumed &
$(1.8637,0.5492)$\\
\bottomrule
\end{tabular}
\end{center}
The self-referential row is a different condition, not a second
solution of the same fixed-input system. Their common physical
prime-lattice origin is not established by this comparison table.

\subsection{Inverse exponent, rounding and back-substitution}
\label{app:inverse-exponent-check}

\begin{proof}[Proof of Proposition~\ref{prop:exponent-5}]
For $g>1$ and $0<\beta<1/\sqrt3$, taking logarithms of the
proposed correction law gives
\[
 \log(\beta\sqrt3)=-\frac{\gamma_1}{4\pi g^n}<0,
 \qquad g^n=\frac{\gamma_1}{4\pi|\log(\beta\sqrt3)|}=:A(g,\beta)>0.
\]
A second logarithm yields $n\log g=\log A(g,\beta)$.
Division by $\log g\ne0$ gives the inverse exponent. This solves
for a calibrated exponent, not for an independently predicted $g$.
\end{proof}

With independent rounded inputs $(g,\beta)=(1.864,0.5493)$,
$A\simeq22.58446676436$ and $n_{\rm exact}\simeq5.00584331607$.
Enforcing $\beta=\sqrt{(g-1)/(g+1)}$ at the same $g$ gives
$\beta\simeq0.549250377928$, $A\simeq22.54357460599$ and
$n_{\rm exact}\simeq5.00293308442$. Both have nearest integer $5$,
but neither establishes a structural integer-exponent law.

As stated explicitly in Remark~\ref{rem:exact-sc}, substitution
$g^{n_{\rm exact}}=\exp(\log A)=A$ is an identity for all allowed
input pairs. It remains an identity after imposing the frame relation
for $1<g<2$. A determining equation requires an independent exponent
law or another independently derived relation.

\subsection{Verification ledger: executed checks and unproved obligations}
\label{app:verification-ledger}

The reproducible algebraic checks are in
\path{scripts/qsc_audit/tests/test_spectral_and_exponent_contracts.py}.
They test exact rational matrix examples, inverse-exponent evaluation,
the derivative-bound factorization and numerical bisection. They do
not prove that an actual MM/HC realization satisfies the hypotheses
or certify the location of a Riemann zero; its ordinate is a declared input.

The relay and readout checks in \path{analysis/OP5_Relay_C5_Check.py}
test declared outgoing models and supplied residual dictionaries.
They do not identify the native C5 residuals, prove an exact
$\Theta$--$R$ balance or establish a gravitational intermediate field.
The framework/adapter checks under \path{scripts/qsc_audit/tests/}
check their implemented grammar and guard behavior, not every formal
MML command or every mathematical assertion of this manuscript.

The executable entry points are
\begin{quote}\small\ttfamily
python -B -m unittest discover -s scripts/qsc\_audit/tests\par
python -B -m unittest discover -s analysis -p OP5\_Relay\_C5\_Check.py
\end{quote}
Reading-route regression tests additionally check unique definition
inclusion and resolution of the new appendix targets. A successful
LaTeX compilation checks document processing, not proof correctness.
No unattended verifier, Monte Carlo balance certificate or automatic
OP closure is claimed behind these references.
